%%%%%%%%%%%%%%%%%%%%%%%
18.
%%%%%%%%%%%%%%%%%%%%%%%
\begin{figure}[h]
	\centering
	\begin{tikzpicture}
		% Axes
		\draw[axis] (-2cm, 0cm) -- (2cm, 0cm) node[right] {$\Re$};
		\draw[axis] (0cm, -2cm) -- (0cm, 2cm) node[right] {$\Im$};

		% Branch cut
		\node[pole] at (0cm, 0cm) {};
		\draw[branch cut] (0cm, 0cm) -- (1.9cm, 0cm);

		% Poles
        	\node[pole] (A) at (0cm, -0.7cm) {}; \node[right] at (A.east) {$-i$};
        	\node[pole] (B) at (0cm,  0.7cm) {}; \node[right] at (B.east) {$i$};

		% Contour
		\def\R{1.8cm}
		\def\eps{0.4cm}
        	%%% TODO: Remove these magical constants
		\draw[contour] (0.7071*\eps, 0.7071*\eps) -- (1.7776cm, 0.7071*\eps) node[below] {$R$};
		\draw[contour] (1.7776cm, 0.7071*\eps) arc (9.0406:350.9093:\R);
		\draw[contour] (1.7776cm, -0.7071*\eps) -- (0.7071*\eps, -0.7071*\eps);
		\draw[contour] (0.7071*\eps, -0.7071*\eps) arc (315:45:\eps);
	\end{tikzpicture}
	\caption{Contour for \refprob{}~3--18.}%
	\label{fig:problem3-18}
\end{figure}

Recalling the exercise in the main text regarding \refeq{3}{37},
we could infer that a direct contour integration of the integrand along some contour analogous to \reffig~3--3
in the textbook would lead to a loss in the exponent of the logarithm (or one could just do it and observe).

If you attempt this, however, one finds that the one needs to know both integrals
\[
	\int_0^\infty dx \frac{\left(\ln x\right)^k}{1+x^2} \quad (k \le 1)
\]
to determine the integral at hand.
Thus, ensuring generality, let
\[
	I_n := \int_0^\infty dx \frac{\left(\ln x\right)^n}{1+x^2}
\]
and $f_n(z) := \frac{{\left( \ln z \right)}^{n+1}}{1 + z^2}$,
where we place the branch cut on the real axis and $\arg z = 0$ just above the real axis.
As such, the poles of $f(z)$ have arguments $\frac{\pi}{2}$ and $\frac{3\pi}{2}$.
Consider the contour shown in \reffig~\ref{fig:problem3-18}.
\begin{align*}
	\oint_C dz f_n(z)
	&= 2i\pi \left( \res{z=i} f_n(z) + \res{z=-i} f_n(z) \right) \\
	&= 2i\pi \left( \evalAt{\frac{{\left( \ln z \right)}^{n+1}}{2z}}{z=i} + \evalAt{\frac{{\left( \ln z \right)}^{n+1}}{2z}}{z=-i} \right) \\
	&= -2i\pi \cdot \frac{\left(3^{n+1}-1\right) i^n \pi^{n+1}}{2^{n+2}}
\end{align*}
We also have
\begin{align*}
    \oint_C dz f_n(z)
    &= \int_\epsilon^R dx \frac{{\left( \ln x \right)}^{n+1}}{1 +x^2}
     + \int_0^{2\pi} d\theta iRe^{i\theta} \frac{{\left( \ln R + i\theta\right)}^{n+1}}{R^2 e^{2i\theta} + 1} \\
    &\quad + \int_R^\epsilon dx \frac{{\left( \ln x + 2i\pi \right)}^{n+1}}{1 + x^2}
     + \int_{2\pi}^0 d\theta i\epsilon e^{i\theta} \frac{{\left(\ln\epsilon + i\theta\right)}^{n+1}}{\epsilon^2 e^{2i\theta} + 1} \\
    &\xrightarrow[\epsilon \rightarrow 0]{R \rightarrow \infty}
        -\sum_{k=0}^n (2i\pi)^{n+1-k} \binom{n+1}{k} I_k
\end{align*}
\[
	\Rightarrow I_n = \inv{n+1} \left(
		\frac{\left(3^{n+1}-1\right) i^n \pi^{n+1}}{2^{n+2}} - \sum_{k=0}^{n-1} (2i\pi)^{n+1-k} \binom{n+1}{k} I_k
	\right)
\]
This recurrence relation, along with $I_0 = \halfpi$, is enough to determine every $I_n$!
\[
	\therefore \int_0^\infty dx \frac{{\ln x}^2}{1 + x^2} = \frac{\pi^3}{8}
\]

Questions remain, however.
For one, it is not obvious from the recurrence relation that $I_n$ must be real.
So, one turns to generating functions for a better understanding of this sequence:
\[
	g(t) := \sum_{n+0}^\infty I_n \frac{t^n}{n!}.
\]
The original integral form lends itself very nicely to cancel the logarithm:
\begin{align*}
	g(t)
	&= \int_0^\infty dx \frac{x^t}{1+x^2} \\
    &= \half \int_0^\infty dx 2x \cdot \left(\frac{x^2}{1+x^2}\right) \cdot \left(x^2\right)^{\frac{t-3}{2}} \\
	&= \int_0^\infty \frac{du}{2(1-u)^2} u \left(\frac{u}{1-u}\right)^{\frac{t-3}{2}} & \left(u := \frac{x^2}{1+x^2}\right) \\
	&= \frac{1}{2} \Beta\left(\frac{1+t}{2}, \frac{1-t}{2}\right) \\
	&= \frac{\Gamma\left(\frac{1+t}{2}\right) \cdot \Gamma\left(\frac{1-t}{2}\right)}{2\Gamma(1)} \\
	&= \frac{\pi \csc\left(\pi\frac{t+1}{2}\right)}{2} & (\Gamma(z)\Gamma(1-z) = \pi\csc\pi z) \\
	&= \halfpi\sec\halfpi t.
\end{align*}
Thus is revealed: the integrals can be succinctly reformulated as
\[
	I_n = \int_0^\infty dx \frac{\left(\ln x\right)^n}{1+x^2}
	= \evalAt{\frac{d^n}{dt^n} \left(\halfpi\sec\halfpi t\right)}{t = 0}
	= \left(\halfpi\right)^{n+1} \evalAt{\frac{d^n}{dt^n} \sec t}{t=0}.
\]
Imagine having to derive this from the recurrence relation, or even going in the reverse direction!

